\section{Bases teóricas}\label{bases-teuxf3ricas}

\subsection{Procesamiento del lenguaje natural}\label{procesamiento-del-lenguaje-natural}

El procesamiento del lenguaje natural (PLN) estudia métodos computacionales para analizar, interpretar y generar lenguaje humano.

En esta investigación se abordará la generación de respuestas textuales breves, no la traducción automática.

\subsection{Arquitectura Transformer}\label{arquitectura-transformer}

La arquitectura Transformer, presentada por \textcite{Vaswani2017}, utiliza mecanismos de atención para representar relaciones entre elementos de una secuencia.

Los modelos de lenguaje autorregresivos basados en esta arquitectura generan texto mediante la predicción sucesiva de tokens.

En la autoatención, las representaciones de consulta ($Q$), clave ($K$) y valor ($V$) determinan cómo se pondera la información de la secuencia. La operación se expresa como:
\begin{equation}
\operatorname{Attention}(Q,K,V)=\operatorname{softmax}\left(\frac{QK^{\mathsf T}}{\sqrt{d_k}}\right)V.
\label{eq:autoatencion}
\end{equation}
En esta expresión, $d_k$ es la dimensión de las claves y el factor de escala evita que los productos internos crezcan con la dimensión.

\subsection{Modelo de lenguaje compacto}\label{modelo-de-lenguaje-compacto}

Para esta investigación, se considerará compacto un modelo con menos de mil millones de parámetros.

Esta definición es operacional y se utilizará únicamente para delimitar el estudio.

El tamaño reducido puede favorecer su utilización en hardware limitado, pero no garantiza competencia en una lengua de pocos recursos.

\subsection{Qwen2.5-0.5B-Instruct}\label{qwen25-05b-instruct}

Qwen2.5-0.5B-Instruct es un modelo autorregresivo ajustado para seguir instrucciones, con aproximadamente 490 millones de parámetros \parencite{QwenCard2024}.

Se utilizará como candidato inicial por su tamaño, disponibilidad pública y licencia Apache 2.0.

Su documentación no demuestra específicamente un rendimiento adecuado en quechua sureño. Por tanto, será necesaria una prueba preliminar de tokenización, generación y compatibilidad técnica.

\subsection{Ajuste fino supervisado}\label{ajuste-fino-supervisado}

El ajuste fino supervisado consiste en adaptar un modelo mediante ejemplos que contienen entradas y salidas esperadas.

En este estudio, las entradas serán preguntas cotidianas y las salidas serán respuestas breves validadas en quechua sureño.

Durante el entrenamiento se procurará que la función de pérdida se aplique principalmente a los tokens de la respuesta, evitando utilizar las instrucciones como objetivo de predicción.

\subsection{QLoRA}\label{qlora}

QLoRA es una técnica de adaptación eficiente que permite entrenar adaptadores de bajo rango sobre un modelo base cuantizado \parencite{Dettmers2023}. Sea $W_0$ la matriz de pesos preentrenada y congelada durante el ajuste; la actualización se parametriza como:

La cuantización reduce el espacio necesario para almacenar determinados pesos, mientras que los adaptadores concentran los parámetros entrenables.

Su aplicación estará condicionada a la compatibilidad del hardware y del software disponible.

En LoRA, la actualización de una matriz de pesos congelada se representa mediante dos matrices entrenables de rango reducido:
\begin{equation}
\Delta W=\frac{\alpha}{r}BA,
\label{eq:descomposicion_lora}
\end{equation}
donde $r$ es el rango y $\alpha$ es el factor de escala. La ecuación describe la parametrización de LoRA; no implica que el ajuste vaya a mejorar la calidad de las respuestas.

\subsection{Generación de respuestas breves}\label{generaciuxf3n-de-respuestas-breves}

La generación de respuestas breves consiste en producir una contestación concisa que atienda una pregunta.

Para este estudio, una respuesta deberá:

\begin{itemize}
\tightlist
\item
  Estar formulada en quechua.
\item
  Corresponder a la variedad solicitada.
\item
  Atender directamente la intención de la pregunta.
\item
  Evitar información no solicitada o inventada.
\item
  Mantener una extensión preferente de una o dos oraciones.
\end{itemize}

La longitud se controlará también mediante un límite técnico de tokens, pero no se equipararán automáticamente los tokens con palabras u oraciones.

\subsection{Quechua sureño}\label{quechua-sureuxf1o}

El estudio comprenderá las variedades chanka y collao.

Se utilizará el \emph{Manual de escritura quechua sureño} del Ministerio de Educación como referencia de escritura, complementado con la revisión de hablantes competentes.

Las decisiones sobre variantes aceptables se documentarán antes de la evaluación.

\subsection{Calidad de respuesta}\label{calidad-de-respuesta}

La calidad de respuesta será un constructo evaluado mediante tres dimensiones:

\textbf{Pertinencia:} correspondencia entre la pregunta y la contestación.

\textbf{Corrección lingüística:} aceptabilidad y comprensibilidad de la formulación.

\textbf{Adecuación a la variedad:} correspondencia con la variedad lingüística solicitada.

Se establecerá un criterio binario de aceptabilidad global a partir de estas dimensiones.

\section{Síntesis del estado del arte y brecha de investigación}\label{suxedntesis-del-estado-del-arte-y-brecha-de-investigaciuxf3n}

La revisión presentada identifica cuatro líneas de investigación relacionadas:

\begin{enumerate}
\def\labelenumi{\arabic{enumi}.}
\tightlist
\item
  Traducción automática quechua--español.
\item
  Construcción y preparación de recursos lingüísticos.
\item
  Adaptación eficiente de modelos.
\item
  Evaluación multilingüe y generación textual.
\end{enumerate}

Los antecedentes demuestran que estas líneas cuentan con investigaciones previas. Por ello, no se planteará como novedad la utilización de QLoRA, la construcción de un modelo de traducción ni la generación de texto en quechua considerada de manera general.

La brecha específica de este estudio se formula de la siguiente manera:

\textbf{Falta establecer, para las condiciones delimitadas de esta investigación, qué efecto produce adaptar mediante QLoRA un único modelo compacto sobre la aceptabilidad de respuestas breves a preguntas cotidianas en quechua sureño, utilizando un corpus validado y una evaluación humana diferenciada para las variedades chanka y collao.}

Esta brecha es una hipótesis de trabajo fundamentada en los antecedentes examinados, no una declaración de inexistencia absoluta de estudios similares.

La contribución esperada será un resultado experimental reproducible, acompañado de un conjunto de evaluación documentado y un análisis de las limitaciones del modelo.

\begin{center}\rule{0.5\linewidth}{0.5pt}\end{center}
